\section{Reading Comprehension Questions}

\noindent
{\em From Section~\ref{section:algorithms}}

\begin{requ}
What does it mean for an algorithm to {\em return a value}?
\begin{requanswer}
It means that the result of the algorithm is some value that can be assigned to a variable.
For instance, when a line of code like \verb+int a = min(4,5)+ executes, 
the \verb+min+ method (algorithm) ``returns'' the minimum value so it can be assigned to 
the variable \verb+a+.
\end{requanswer}
\end{requ}

\begin{requ}
Write a method that computes the volume of a cuboid 
(or rectangular prism)
with width $w$, height $h$ and depth $d$, where each is a real number.
\begin{requanswer}
Here it is. Note that the order of the parameters does not matter.
\begin{code}
double cuboidVolume(double w, double h, double d) {
    return w*h*d;
}
\end{code}
\end{requanswer}
\end{requ}


\medskip 

\noindent
{\em From Section~\ref{section:moddivRepeat}}

\begin{requ}
Write a method \verb+modN(int a, int n)+ that computes $a\pmod n$ 
{\em without using the modulus operator}. That is, you may only use basic integer arithmetic
(e.g. $+$, $-$, $*$, $/$).
Your code should be as simple and efficient as possible.
You may assume that $a\geq 0$ if it helps.
(Hint: Assume that integer division truncates.) 
\begin{requanswer}
Here is one possibility: keep subtracting $n$ from $a$
until you get to a number less than $n$:
\begin{code}
int modN(int a, int n) {
    while(a>=n) {
        a=a-n;
    return a;
}
\end{code}
That is not very efficient if $a$ is large, however. 
A more efficient solution would be as follows:
\begin{code}
int modN(int a, int n) {
    return a - (a/n)*n;
}
\end{code}
Make sure you understand why this solution works!
Also, it should be noted that both of these solutions require that $a\geq 0$.
\end{requanswer}
\end{requ}

\begin{requ}
If you did not get Exercise~\ref{exer:roundNotTruncatetwo}
correct the first time,
try doing it again without looking at the answer first!
\begin{requanswer}
As was the case there,
two reasonable solutions include \verb@(n+m/2)/m@ and \verb@(2n+m)/(2m)@.
\end{requanswer}
\end{requ}


\medskip 

\noindent
{\em From Section~\ref{section:if}}

\begin{requ}
Write a method \verb+boolean congruentModN(int a, int b,int n)+
that returns true if and only if $a\equiv b \pmod n$.
Your code should be as simple and efficient as possible.
\begin{requanswer}
Here is one possible solution:
\begin{code}
boolean congruentModN(int a, int b,int n) {
    if(a%n==b%n) {
        return true;
    } else {
        return false;
    }
}
\end{code}
Here is a very concise solution. Do you see why it works?
\begin{code}
boolean congruentModN(int a, int b,int n) {
    return((a-b)%n == 0);
}
\end{code}
\end{requanswer}
\end{requ}

\begin{requ}
Does your solution to the previous question still work if 
the language you are working with has a mod operator
that can return a negative value? If so, fix it.
\begin{requanswer}
The first one will {\em not} work if $a$ and $b$ have different signs.
The second one will always work since if the difference between
$a$ and $b$ is a multiple of $n$, it will always return $0$.
Thus, the second solution is better.
\end{requanswer}
\end{requ}


\begin{requ}
Write a conditional statement that determines if $x$ is odd. 
Your statement should be as short as possible.
\begin{requanswer}
This one is pretty straightforward: \verb+if(x%2==1)+
\end{requanswer}
\end{requ}

\begin{requ}
Write a segment of code that adds 1 to $x$ if $x$ is even and subtracts 1 from $x$ if $x$ is odd.
For an added challenge, see if you can write a second version that does not use a conditional statement.
\begin{requanswer}
The easy answer is as follows:
\begin{code}
if(x%2==0) {
   x++; 
} else {
    x--;
}
\end{code}
One way to do it without a conditional statement is: \verb^x+=1-2*(x%2);^ Yikes!
\end{requanswer}
\end{requ}

\medskip

\noindent
{\em From Section~\ref{section:logicInProgramming}}
\begin{requ}
Consider the following code: \verb+if(...) { // do something }+ 
(where the details of the conditional statement in the \verb+if+ have been omitted).
Is it more likely that the conditional statement is a tautology or a contingency? Explain.
\begin{requanswer}
It is probably a contingency. 
A tautology is a proposition that is always true. Since there is a conditional statement involved, it is likely that different cases happen.
In other words, there would be no need for the \verb+if+ statement if the expression inside it is a tautology.
\end{requanswer}
\end{requ}

\begin{requ}
If a programming language does not implement short-circuiting with logical expressions, how
would you write the following code instead so that it does not 
ever crash? 
\begin{code}
    if( !list.isEmpty() && list.get(0)>= 100 ) {
        list.clear();
    }
\end{code}
\begin{requanswer}
We need to use a nested conditional statement:
\begin{code}
if(!list.isEmpty()) {
    if(list.get(0)>= 100) {
        list.clear(); 
    }
}
\end{code}
\end{requanswer}
\end{requ}

%% Moved to later section. Keeping here to make sure 
%% it didn't get deleted accidentally.
%\begin{requ}
%Consider the conditional statements \verb+if( i<a.length && a[i]!=0 )+
%and\\
% \verb+if( a[i]!=0 &&  i<a.length )+. Are they equivalent in most programming languages?
% Explain why or why not.
%If they are not equivalent, explain exactly how they differ. In particular, is one right and one
%wrong, or do they accomplish different goals?
%\begin{requanswer}
%Although {\em logically} these two conditional statements are equivalent, they are actually not equivalent
%in most programming languages due to short-circuiting. The first one is correct because it makes sure
%that the index is valid before indexing into the array. 
%The second one will probably crash if $i$ indexes outside of the array, so it is not correct.
%\end{requanswer}
%\end{requ}

\begin{requ}
Simplify the conditional statement \verb+if( !(x<=0 || y<=0) )+ as much as possible.
\begin{requanswer}
\verb+if( x > 0 && y > 0 )+ (using DeMorgan's law). Note that 
\verb+if( x >= 0 && y >= 0 )+ is {\bf not} correct since \verb+x>0+ is the negation of \verb+x<=0+.
\end{requanswer}
\end{requ}


\medskip

\newpage

\noindent
{\em From Section~\ref{section:bit}}
\begin{requ}
Assume $x$ and $y$ are integers stored on a computer using 8 bits.
Let $x=37$ ($00100101$ in binary), $y=112$ ($01110000$ in binary).
Compute \verb+~+$x$, $x \verb+&+ y$, $x \verb+|+ y$, and $x \verb+^+ y$. Give your answers in both binary and decimal.
\begin{requanswer}
$\verb+~+x = 218$   (11011010), 
$x \verb+&+ y = 32$ (00100000),
$x \verb+|+ y = 117$ (01110101), and
$x \verb+^+ y = 85$ (01010101).
\end{requanswer}
\end{requ}

\medskip

\noindent
{\em From Section~\ref{section:for}}

\begin{requ}\label{simplefor}
Exactly how many times does the following loop execute?
\begin{code}
  for(int i=1;i<n;i++) {
     // do something
   }
\end{code}
\begin{requanswer}
$n-1$. Notice that the loop starts at 1 and ends at $n-1$, so that's $n-1$ times 
(not $n$ as you might have guessed).
\end{requanswer}
\end{requ}

\begin{requ}
Write a method \verb+sumFirstN(int n)+ that returns the sum of the first $n$ positive integers.
Your code should be as simple and efficient as possible.
(Hint: you should use a loop, although we will see later on this semester how to solve this particular
problem without a loop.)
\begin{requanswer}
Here is the solution I expected you to come up with:
\begin{code}
sumFirstN(int n) {
	int sum = 0; 
	for( int i = 1; i <= n; i++) {
		sum+=i; 
	}
}
\end{code}
O.K., I can't resist giving the better answer. The following is much more efficient, but it may
be unclear to you why it is correct. Do not worry--we will see the idea behind it later.
\begin{code}
sumFirstN(int n) {
	return n*(n + 1) / 2;
}
\end{code}
Do not worry if you do not understand this one---I 
definitely did {\em not} expect you to come up with this one yet.
But I wanted to prime your mind for something we will learn very soon.
\end{requanswer}
\end{requ}

\medskip 

\noindent
{\em From Section~\ref{section:arrays}}

\begin{requ}
Write a method called \verb+int minimum(int a[], int n)+, where $a$ is an array of length $n$,
that returns the smallest element in the array. 
Your code should be as simple and efficient as possible.
(Note: You {\em cannot} assume that the entries of the array are all positive.)
\begin{requanswer}
\begin{code}
int minimum(int a[], int n) {
    int min = a[0];
    for(int i=1;i<n;i++) {
        if(a[i]<min) {
            min=a[i];
        }
    }
    return min;
}
\end{code}
\end{requanswer}
\end{requ}

\begin{requ}
Write a method called \verb+boolean containsZero(int a[], int n)+, where $a$ is an array of length $n$,
that returns true if and only if some element of $a$ is 0 (that is, \verb+a[i]==0+ for some $0\leq i<n$). 
Your code should be as simple and efficient as possible.
\begin{requanswer}
Does your solution use a \verb+boolean+ variable to keep track of whether or not you found a 0? 
If so, try to rewrite your code so that it does not do so before you read the solution given below.
If you used a boolean variable, your solution is both more complicated and less efficient than it should be,
and it is best that you figure out the better way on your own so that you are less likely to write similar
code in the future.

Again, before reading the solution given below, does your solution contain an \verb+else+ clause? 
If so, think about whether or not it is correct by walking through your code on a small array.
Since I can't read your code, I cannot be certain, but if you have an \verb+else+ clause, 
your solution is likely to be incorrect.
\begin{code}
boolean containsZero(int a[], int n) {
    for(int i=0;i<n;i++) {
        if(a[i]==0) {
            return true;
        }
    }
    return false;
}
\end{code}
\end{requanswer}
\end{requ}

\begin{requ}\label{exer:noZeroes}
Write two versions of a method called \verb+boolean noZeroes(int a[], int n)+, 
where $a$ is an array of length $n$,
that returns true if and only if none of the elements of $a$ are 0: 
one that calls \verb+containsZero+ (see the previous question), and one that does not.
In both cases, your code should be as simple and efficient as possible.
\begin{requanswer}
Here is the first version:
\begin{code}
boolean noZeroes(int a[], int n) {
    return !containsZero(a,n);
}
\end{code}
Here is the second version, which looks almost identical to the solution to the previous question:
\begin{code}
boolean noZeroes(int a[], int n) {
    for(int i=0;i<n;i++) {
        if(a[i]==0) {
            return false;
        }
    }
    return true;
}
\end{code}
\end{requanswer}
\end{requ}


\begin{requ}
Consider the conditional statements 
\verb+if(i>=0 && i<a.length && a[i]!=0)+
and
 \verb+if(a[i]!=0 &&  i<a.length && i>=0)+. Are they equivalent in most programming languages?
 Explain why or why not.
If they are not equivalent, explain exactly how they differ. In particular, is one right and one
wrong, or do they accomplish different goals?
\begin{requanswer}
Although these two conditional statements are 
{\em logically equivalent}, they are actually not equivalent
in most programming languages due to short-circuiting. 
The first one is correct because it makes sure
that the index is valid before indexing into the array. 
The second one will probably crash if $i$ indexes outside of the array, so it is not correct.
For a more complete discussion, see Example~\ref{exa:ss}, since
it is addresses almost the same exact same question! 
The only difference is that here we added a check to make sure the
index was not negative.
\end{requanswer}
\end{requ}

\medskip

\noindent
{\em From Section~\ref{section:quantsAndAlgs}}

\begin{requ}\label{rq:fooBeforeBar}
Let $A$ be an array. Consider the statement $\forall x (A[x]!=0)$,
where the universe of discourse is $\{0,1,\ldots, n-1\}$.
\begin{enumerate}[(a)]
\item
Rephrase the statement in English.
\item
Write a function \verb+boolean foo(int []A, int n)+ 
that computes and returns the truth value of the statement.
\item
What is a better name for \verb+foo+? In other words, what does it do?
\end{enumerate}
\begin{requanswer}
(a) The array has no entries that are 0.
(b) Here is the most obvious solution:
\begin{code}
boolean foo(int []A, int n)
    for (int i = 0; i < n; i++) {
        if(a[i] == 0) {
            return false;  
        }
    }
    return true; 
}
\end{code}
(c) A better name would be \verb+noZeroes+ or \verb+containsNoZeroes+
because it determine whether or not the array contains any zeroes, returning true if it contains no zeroes and 
false otherwise.
\end{requanswer}
\end{requ}


\begin{requ}
Let $A$ be an array. Consider the statement $\neg \exists x (A[x]==0)$,
where the universe of discourse is $\{0,1,\ldots, n-1\}$.
\begin{enumerate}[(a)]
\item
Rephrase the statement in English.
\item
Write a function \verb+boolean bar(int []A, int n)+ 
that computes and returns the truth value of the statement.
\item
What is a better name for \verb+bar+? In other words, what does it do?
\end{enumerate}
\begin{requanswer}
(a) Translating it very literally, it is saying that it is not the case that there exists an element of the array that is zero.
In other words, the array contains no zeroes. In other words, this is saying the exact same thing as the expression 
in Question~\ref{rq:fooBeforeBar}.
(b) The same algorithm from Question~\ref{rq:fooBeforeBar}, but replace \verb+foo+ with \verb+bar+.
(c) Same as answer to Question~\ref{rq:fooBeforeBar}.
\end{requanswer}
\end{requ}

\begin{requ}
Given an array $a$ of length $n$, give an algorithm that will determine the truth value of
$\forall x ( (a[x]\mbox{ is even }) \vee (a[x] \mbox{ is a multiple of 3}) )$.
\begin{requanswer}
\begin{code}
boolean allAre2MultipleOr3Multiple(int[] a, int n) {
   for(int i=0;i<n;i++) {
       if(a[i]%2!=0 && a[i]%3!=0) {
           return false;
       }
   }
   return true;
}
\end{code}
\end{requanswer}
\end{requ}

\medskip 

\noindent
{\em From Section~\ref{section:while}}
\begin{requ}
Rewrite the code from Questions~\ref{simplefor} using a {\em while} loop.
\begin{requanswer}
\begin{code}
int i=1;
while(i<n) {
    // do something
    i++;
}
\end{code}
\end{requanswer}
\end{requ}

\begin{requ}
Rewrite the code from Questions~\ref{exer:noZeroes} using 
a {\em while} loop that does not contain a return statement
in the loop.
\begin{requanswer}
Go through each element until you find a zero.
At the end, if we looked at the last element, there are
no zeros, so we return true. Otherwise, we return false.
\begin{code}
boolean noZeroes(int a[], int n) {
    int i=0;
    while(i<n && a[i]!=0) {
        i++;
    }
    if(i==n) {
        return true;
    } else {
        return false;
    }
}
\end{code}
\end{requanswer}
\end{requ}



